\documentclass[10pt,a4paper]{amsart}
\usepackage[margin=19mm]{geometry}
\usepackage{amsmath,amssymb,mathtools}
\usepackage[scr=rsfs,cal=euler]{mathalfa}
\usepackage{xfrac}
\usepackage[unicode,hidelinks,bookmarks=false]{hyperref}
\newcommand{\ie}{i.e.\ }
\newcommand{\cf}{cf.\ }
\newcommand{\resp}{respectively\ }
\newcommand{\Subsection}{\subsection}
\providecommand{\nobreakdash}{\nobreak\hbox{-}}
\setlength{\emergencystretch}{2em}
\multlinegap0pt
\usepackage{bm}
\usepackage{bbm}
\usepackage{dsfont}
\usepackage{xspace}
\usepackage{cancel}
\usepackage{mathtools}
\usepackage{amsthm}
\makeatletter
\@ifundefined{lemma}{\newtheorem{lemma}{Lemma}}{}
\makeatother
\DeclareMathOperator{\Bernoulli}{Bern}
\DeclareMathOperator{\Binomial}{Binom}
\newcommand{\EE}{{\mathbb{E}}}
\title[On Unified and Sharpened CMI Bounds for Generalization Error]{On Unified and Sharpened CMI Bounds for Generalization Errors}
\author{Yang Lu, Matthias Frey, Margreta Kuijper, Jingge Zhu}
\date{Source: arXiv:2605.21056v1 (CC BY 4.0)}
\begin{document}
\maketitle

\noindent\emph{Source and licence.} On Unified and Sharpened CMI Bounds for Generalization Errors. Version arXiv:2605.21056v1 (CC BY 4.0), distributed under the Creative Commons Attribution 4.0 International licence, as recorded in the arXiv licence metadata for this exact version and shown on the abstract page https://arxiv.org/abs/2605.21056v1. Full rights notice: licenses/arxiv-2605.21056-CC-BY-4.0.txt.\par\smallskip

\noindent\textit{Editorial scope.} Counted unit: the Lemma whose source label is \texttt{lem:exp\_bin\_logx} in the article (source0.tex lines 1846--1853, source label \texttt{lem:exp\_bin\_logx}), delivered together with the article's own complete original proof of it (source0.tex lines 1855--1923, one \texttt{proof} environment of 69 lines). No mathematics is written, repaired or re-derived: statement and proof are the article's own text line for line. Required context retained uncounted: none. Every definition, display and label of those lines is retained unchanged. Omitted from this package and not redistributed: 1--1845, 1854--1854, 1924--2706. Nothing inside a retained excerpt is omitted or altered: every retained body is delivered exactly as the article's lines are written, so no omission receipt of any kind accompanies this package. Citations: the retained excerpts carry no citation command at all, so no bibliography entry of the article is transcribed and no reference is redistributed. Reference adaptation: none was needed and none was made; every reference inside the delivered unit resolves inside it, so no unresolved reference or citation is produced. Printed slips kept literal: no printed word, symbol, hypothesis or number of the counted statement or of its proof was altered. Layout and class: the article's own class is \texttt{IEEEtran} with packages including cite, amsmath, amssymb, amsfonts, amsthm, mathrsfs, bm, dsfont, algorithm2e, mathtools, graphicx, textcomp, multirow, threeparttable; the wrapper is the lane's pinned \texttt{amsart} editorial wrapper at 10pt on A4, which restates the theorem-like environments the retained text needs and adds the article's own macro definitions that the retained text needs (\texttt{\textbackslash Bernoulli}, \texttt{\textbackslash Binomial}, \texttt{\textbackslash EE}), with extra wrapper packages (bm, bbm, dsfont, xspace, cancel, mathtools). The delivered numbering is the unit's own. Assets: no retained span contains a figure environment, a graphics-inclusion command, a PDF-page inclusion, a table or a TikZ picture; no image file is copied into this package, the package declares no asset dependency and no image file is redistributed with it. Licence: version arXiv:2605.21056v1 of this article is distributed under the Creative Commons Attribution 4.0 International licence, as recorded in the arXiv licence metadata for this exact version and shown on the abstract page https://arxiv.org/abs/2605.21056v1. A full rights notice is delivered at licenses/arxiv-2605.21056-CC-BY-4.0.txt. Characters: every retained excerpt is pure ASCII, so the delivered engine, XeLaTeX with its default Latin Modern text font, reports no missing character.
	\begin{lemma} \label{lem:exp_bin_logx}
		If $X$ is a binomial random variable following $X \sim \Binomial(n, p)$ and $a \geq 0$, then
		\begin{equation}
			\begin{aligned}
				\EE \left[ \log (X + a) \right] = \log (np + a) - \frac{np (1 - p)}{2 (np + a)^2} + O\left( \frac{1}{n^2} \right).
			\end{aligned}
		\end{equation}
	\end{lemma}

	\begin{proof}
		By Taylor's theorem, we expand the function $f(x) = \log (x + a)$ at point $x_0$ such that
		\begin{equation}
			\begin{aligned}
				f(x) = f(x_0) + \sum_{i=1}^{3} \frac{f^{(i)}(x)}{i !} (x - x_0)^i + \frac{f^{(i)}(\xi)}{4 !} (x - x_0)^4,
			\end{aligned}
		\end{equation}
		where $\xi$ is between $x_0$ and $x$. Let $x_0 = np$, it can be further calculated that
		\begin{equation}
			\begin{aligned}
				\log (x + a) =& \log (np + a) + \frac{x - np}{np + a} - \frac{(x - np)^2}{2 (np + a)^2} + \frac{(x - np)^3}{3 (np + a)^3} - \frac{(x - np)^4}{4 (\xi + a)^4}.
			\end{aligned}
			\label{eq:taylor_logx}
		\end{equation}
		Taking expectation on the upper bound of \eqref{eq:taylor_logx} yields
		\begin{equation}
			\begin{aligned}
				\EE \left[ \log (X + a) \right] &= \EE \left[ \log (np + a) + \frac{X - np}{np + a} - \frac{(X - np)^2}{2 (np + a)^2} + \frac{(X - np)^3}{3 (np + a)^3} \right] - \EE \left[ \frac{(X - np)^4}{4 (\xi + a)^4} \right] \\
				&= \log (np + a) + \frac{\EE \left[ X - np \right]}{np + a} - \frac{\EE \left[ (X - np)^2 \right]}{2 (np + a)^2} + \frac{\EE \left[ (X - np)^3 \right]}{3 (np + a)^3} - \EE \left[ \frac{(X - np)^4}{4 (\xi + a)^4} \right] \\
				&= \log (np + a) + 0 - \frac{np (1 - p)}{2 (np + a)^2} + \frac{np (1 - p) (1 - 2p)}{3 (np + a)^3} - \EE \left[ \frac{(X - np)^4}{4 (\xi + a)^4} \right] \\
				&= \log (np + a) - \frac{np (1 - p)}{2 (np + a)^2} + O\left( \frac{1}{n^2} \right) - \EE \left[ \frac{(X - np)^4}{4 (\xi + a)^4} \right].
			\end{aligned}
		\end{equation}
		It suffices to show that the expectation of the remainder term is of order $O(\frac{1}{n^2})$. Using Hoeffding's inequality, we can derive an upper bound for the tail probability that $X \leq \frac{1}{2} np$, which is shown as follows:
		\begin{equation}
			\begin{aligned}
				\Pr\left[ X \leq \frac{1}{2} np \right] &= \Pr\left[ X - \EE \left[ X \right] \leq -\frac{1}{2} np \right] \\
				&= \Pr\left[\sum_{i=1}^{n} X_i - \sum_{i=1}^{n} \EE \left[ X_i \right] \leq - \frac{1}{2} np \right] \\
				&\leq e^{-\frac{1}{2} p^2 n},
			\end{aligned}
		\end{equation}
		where $X_i \sim \Bernoulli(p)$ is a Bernoulli random variable for any $i \in [1, n]$. Then, by the tower property of expectation, we rewrite the expectation of the remainder term as
		\begin{equation}
			\begin{aligned}
				\EE \left[ \frac{(X - np)^4}{4 (\xi + a)^4} \right] &= \Pr\left[ X > \frac{1}{2} np \right] \EE \left[ \left. \frac{(X - np)^4}{4 (\xi + a)^4} \right| X > \frac{1}{2}np \right] + \Pr\left[ X \leq \frac{1}{2} np \right] \EE \left[ \left. \frac{(X - np)^4}{4 (\xi + a)^4} \right| X \leq \frac{1}{2}np \right].
			\end{aligned}\label{eq:exp_remainder}
		\end{equation}
		Since $X \in [n]$ and $\xi$ is within $np$ and $X$, we have 
		\begin{equation}
			\begin{aligned}
				\frac{(X - np)^4}{4 (\xi + a)^4} \leq \frac{\max((np-1)^4, (n-np)^4)}{4(1+a)^4} = O(n^4).
			\end{aligned}
		\end{equation}
		Due to fact that exponential decay dominates any polynomial growth, we see that the contribution of $X \leq \frac{1}{2} np$ to the expectation is negligible, i.e.,
		\begin{equation}
			\begin{aligned}
				\Pr\left[ X \leq \frac{1}{2} np \right] \EE \left[ \left. \frac{(X - np)^4}{4 (\xi + a)^4} \right| X \leq \frac{1}{2}np \right] \leq e^{-\frac{1}{2} p^2 n} \frac{\max((np-1)^4, (n-np)^4)}{4(1+a)^4} = o(\frac{1}{n^2}).
			\end{aligned}\label{eq:exp_X_out_range}
		\end{equation}
		In addition, similar to \eqref{eq:exp_remainder}, we have
		\begin{equation}
			\begin{aligned}
				\EE \left[ (X - np)^4 \right] &= \Pr\left[ X > \frac{1}{2} np \right] \EE \left[ \left. (X - np)^4 \right| X > \frac{1}{2}np \right] + \Pr\left[ X \leq \frac{1}{2} np \right] \EE \left[ \left. (X - np)^4 \right| X \leq \frac{1}{2}np \right] \\
				&\geq \Pr\left[ X > \frac{1}{2} np \right] \EE \left[ \left. (X - np)^4 \right| X > \frac{1}{2}np \right] \\
				&\geq \left( 1 - e^{-\frac{1}{2} p^2 n} \right) \EE \left[ \left. (X - np)^4 \right| X > \frac{1}{2}np \right].
			\end{aligned}\label{eq:exp_X_in_range}
		\end{equation}
		Now, by using the results in \eqref{eq:exp_X_out_range} and \eqref{eq:exp_X_in_range}, we can bound \eqref{eq:exp_remainder} as
		\begin{equation}
			\begin{aligned}
				0 \leq \EE \left[ \frac{(X - np)^4}{4 (\xi + a)^4} \right] &\overset{\eqref{eq:exp_X_out_range}}{\leq} \EE \left[ \left. \frac{(X - np)^4}{4 (\xi + a)^4} \right| X > \frac{1}{2}np \right] + o(\frac{1}{n^2}) \\
				&\leq \EE \left[ \left. \frac{(X - np)^4}{4 (\frac{1}{2}np + a)^4} \right| X > \frac{1}{2}np \right] + o(\frac{1}{n^2}) \\
				&\overset{\eqref{eq:exp_X_in_range}}{\leq} \frac{1}{1 - e^{-\frac{1}{2} p^2 n}} \frac{1}{4 (\frac{1}{2}np + a)^4} \EE \left[ (X - np)^4 \right] + o(\frac{1}{n^2}) \\
				&= \frac{1}{1 - e^{-\frac{1}{2} p^2 n}} \frac{np(1-p)(1+ 3(n-2)p(1-p))}{4 (\frac{1}{2}np + a)^4} + o(\frac{1}{n^2}) \\
				&= \Theta(\frac{1}{n^2}),
			\end{aligned}
		\end{equation}
		which implies \eqref{eq:exp_remainder} is of order $O(\frac{1}{n^2})$. The proof is complete.
	\end{proof}

\end{document}
